\section{Background}

\subsection{Paper 10's Mixed Outcome}

Paper~10~\cite{paper10adaptive} introduced a one-threshold magnitude
test as the simplest deployable change-point detector for online
recalibration. At pre-registered $\tau = 0.05$ the procedure produced
H1 supported (no hazard at $K = 200$, cell-mean delta $+0.7$~pp
vs fixed), H2 rejected ($K = 50$ adaptive falls $5.3$~pp below
lag-1), H3 rejected (fire fractions non-monotone in $K$), and H4
supported (adaptive marginally beats static gate by $+0.7$~pp).
Paper~10's discussion (\S8.3) argued the H1 success and H2 cost are
not separable: any non-trivial adaptive procedure must, by definition,
sometimes gate out gainful calibration, and the gating rate is
parameterised by $\tau$.

\subsection{The Hypothesised Trade-Off}

The intuition Paper~10 left unverified is: smaller $\tau$ should
fire less often, shrinking the H2 penalty at moderate capacity but
re-introducing some H1 hazard at high capacity; larger $\tau$ should
fire more often, eliminating H1 hazard but enlarging the H2 cost.
The relationship should produce a Pareto curve over which an operator
can choose an operating point.

\subsection{Why a Sensitivity Sweep First}

Before evaluating more sophisticated detectors (CUSUM, Bayesian),
it is useful to know whether the simple magnitude detector --- viewed
as a one-parameter family --- contains an operating point that
satisfies \textit{both} Paper~10 pre-registered tolerances. If yes,
the operating point is the deployable recipe. If no, the simple
detector is fundamentally insufficient and a richer detector class
is necessary.
